\chapter{Navigieren, Auswählen, Ein- und Ausblenden}\label{kap:navigation}
\index{Navigation}

Dieses Kapitel beschreibt, wie Sie sich in den vier Ansichten bewegen, wie Sie Objekte und
Teilobjekte auswählen und wie Sie Objekte vorübergehend ausblenden. Diese Handgriffe brauchen
Sie bei jeder Arbeit; es lohnt sich, sie an einem Beispielmodell auszuprobieren:
Öffnen Sie dazu \datei{examples/demo-parts.step} (\autoref{sec:beispiele}).

\section{Die aktive Ansicht}\label{sec:aktive-ansicht}

Jede Mausaktion wirkt in der Ansicht, in der sie stattfindet. Ein Klick (mit einer
beliebigen Maustaste) macht eine Ansicht außerdem zur \emph{aktiven Ansicht}: Sie erhält einen
hervorgehobenen Rahmen, und die Statusleiste zeigt \enquote{Active view: \dots}. Auf die aktive
Ansicht wirken

\begin{itemize}
  \item die Darstellungsmodi aus Werkzeugleiste und Menü \menu{View},
  \item \gui{Zoom extents} und \gui{Zoom selected} (Knopf, Menü und Tastenfolge),
  \item das Maximieren mit der \keys{Leertaste}.
\end{itemize}

Mit gedrückter \keys{Shift}-Taste wirken die Knöpfe der Werkzeugleiste stattdessen auf alle
vier Ansichten.

\section{Drehen, Verschieben, Zoomen}\label{sec:orbit-pan-zoom}
\index{Orbit}\index{Pan}\index{Zoom}

\begin{referenz}{>{\raggedright\arraybackslash}p{3.4cm}>{\raggedright\arraybackslash}p{5.2cm}X}
\toprule
\textbf{Aktion} & \textbf{Maus / Tasten} & \textbf{Hinweise} \\
\midrule
\endhead
Drehen (Orbit) & \maus{rechts ziehen} oder \keys{Alt}+\maus{links ziehen} & nur in der
  perspektivischen Ansicht; in den anderen Ansichten verschieben diese Gesten. \\
Verschieben (Pan) & \maus{mittlere Taste ziehen}, \keys{Shift}+\maus{rechts ziehen} & in
  allen Ansichten; in den Parallelansichten auch \maus{rechts ziehen} und
  \keys{Alt}+\maus{links ziehen}. \\
Zoomen am Cursor & \maus{Mausrad} & Der Punkt unter dem Mauszeiger bleibt stehen. Eine Raste
  ändert den Maßstab um etwa \SI{15}{\percent}. \\
Alles zeigen (Zoom extents) & \keys{Z}, dann \keys{E}; \maus{Doppelklick mit der mittleren
  Taste}; Knopf \gui{Zoom extents} & Doppelklick wirkt auf die angeklickte Ansicht, Taste und
  Knopf auf die aktive; \keys{Shift}+Knopf auf alle. \\
Auswahl zeigen (Zoom selected) & \keys{Z}, dann \keys{S}; Knopf \gui{Zoom selected} & ohne
  Auswahl wie Zoom extents. \\
Maximieren / zurück & \keys{Leertaste} & aktive Ansicht. \\
Befehl beenden / wiederholen & \maus{Rechtsklick ohne Ziehen} & wie \keys{Enter}: beendet den
  laufenden Befehl, ohne laufenden Befehl startet er den letzten erneut. Nur ein Ziehen dreht
  bzw.\ verschiebt. \\
\bottomrule
\end{referenz}

\paragraph{Drehen.} Gedreht wird wie auf einem Drehteller: Seitliches Ziehen dreht um die
senkrechte $Z$-Achse, Ziehen nach oben oder unten kippt die Blickrichtung. Die Kamera kippt
dabei nie über den Pol, sodass $Z$ immer nach oben zeigt. Drehpunkt ist die Mitte der Ansicht;
verschieben Sie das Modell also zuerst so, dass der interessierende Teil in der Mitte liegt,
oder benutzen Sie \gui{Zoom selected}, das die Auswahl in die Mitte rückt.

\paragraph{Verschieben.} Das Modell folgt dem Mauszeiger genau, unabhängig vom Zoom. Während
des Ziehens zeigt der Mauszeiger eine geschlossene Hand.

\paragraph{Zoomen.} Das Mausrad vergrößert bzw.\ verkleinert um den Punkt unter dem
Mauszeiger. Zeigen Sie also auf das Detail, das Sie vergrößern wollen, und drehen Sie das
Rad.

\begin{tipp}
Wer viel mit der Maus arbeitet: \maus{mittlere Taste} zum Verschieben, \maus{rechte Taste}
zum Drehen, \maus{Mausrad} zum Zoomen, \maus{Doppelklick mittlere Taste} für alles. Mit einem
Mac-Trackpad sind die Gesten des nächsten Abschnitts bequemer.
\end{tipp}

\subsection{Trackpad}\label{sec:trackpad}
\index{Trackpad}

Auf dem Mac versteht das Programm die üblichen Trackpad-Gesten:

\begin{referenz}{>{\raggedright\arraybackslash}p{6.0cm}X}
\toprule
\textbf{Geste} & \textbf{Wirkung} \\
\midrule
\endhead
mit zwei Fingern wischen & verschieben \\
zwei Finger auseinander- oder zusammenziehen (Pinch) & zoomen am Cursor \\
\keys{Cmd} + mit zwei Fingern wischen & zoomen am Cursor \\
mit zwei Fingern doppelt tippen & alles zeigen (Zoom extents) \\
\keys{Alt} + klicken und ziehen & drehen (perspektivische Ansicht) \\
\keys{Shift} + Sekundärklick und ziehen & verschieben \\
\bottomrule
\end{referenz}

Unter Windows und Linux hängt das Verhalten eines Touchpads vom Treiber ab: Meldet er das
Wischen mit zwei Fingern als Geste, verschiebt es die Ansicht wie auf dem Mac, sonst zoomt es
wie ein Mausrad; mit \keys{Strg} zoomt es immer.

\subsection{Alles zeigen: Zoom extents}\label{sec:zoom-extents}
\index{Zoom extents}

\gui{Zoom extents} stellt den Zoom so ein, dass alle sichtbaren Objekte in die Ansicht passen;
ausgeblendete zählen nicht mit. Ist nichts sichtbar, zeigt die Ansicht einen Bereich von
rund $\pm\SI{10}{mm}$ um den Ursprung. Nach dem
Öffnen einer Datei und nach dem Ende einer Plugin-Rechnung, die ein leeres Dokument gefüllt
hat, führt das Programm \gui{Zoom extents} in allen Ansichten selbst aus.

\gui{Zoom selected} tut dasselbe für die Auswahl -- ganze Objekte, aber auch einzelne Flächen,
Kanten oder Eckpunkte (\autoref{sec:teilobjekte}). Für einen einzelnen Punkt oder ein sehr
kleines Element zeigt die Ansicht eine Umgebung von etwa \SI{2}{\percent} der Modellgröße um
das Element. Ohne Auswahl wirkt \gui{Zoom selected} wie \gui{Zoom extents}.

\begin{hinweis}
Die Tastenfolgen \keys{Z}, \keys{E} und \keys{Z}, \keys{S} drücken Sie nacheinander, nicht
gleichzeitig: erst \keys{Z} loslassen, dann \keys{E} bzw.\ \keys{S}.
\end{hinweis}

\subsection{Eine Ansicht maximieren}\label{sec:maximieren}
\index{Maximieren}

Die \keys{Leertaste} vergrößert die aktive Ansicht auf die ganze Fläche der vier Ansichten;
ein zweiter Druck stellt das Raster $2 \times 2$ wieder her. Dasselbe bewirkt der Eintrag
\gui{Maximize / restore (Space)} im Ansichtsmenü. Klicken Sie also zuerst in die Ansicht, die
Sie groß sehen wollen, und drücken Sie dann die \keys{Leertaste}.

\subsection{Das Ansichtsmenü}\label{sec:ansichtsmenue}
\index{Ansichtsmenü}

Ein Klick mit der linken Maustaste auf die Kopfzeile einer Ansicht (\enquote{Front (XZ)
$\cdot$ Shaded $\blacktriangledown$}) öffnet ihr Ansichtsmenü:

\begin{referenz}{>{\raggedright\arraybackslash}p{4.8cm}X}
\toprule
\textbf{Eintrag} & \textbf{Wirkung} \\
\midrule
\endhead
\gui{Front (XZ)}, \gui{Left (YZ)}, \gui{Top (XY)}, \gui{Perspective} & Macht aus dieser
  Ansicht die gewählte Standardansicht, mit deren Blickrichtung, Projektion und
  Konstruktionsebene. Die aktuelle ist abgehakt. Der Zoom wird dabei auf die Umgebung des
  Ursprungs zurückgesetzt; wählen Sie danach \gui{Zoom extents}. \\
\gui{Perspective projection} & Schaltet zwischen perspektivischer Projektion (Häkchen) und
  Parallelprojektion um. So lässt sich z.\,B.\ die perspektivische Ansicht parallel
  darstellen, um Maße ohne Verzerrung zu vergleichen. Die Kopfzeile zeigt die Abweichung vom
  Standard als \enquote{(parallel)} bzw.\ \enquote{(perspective)}. \\
\gui{Zoom extents} & ganzes Modell in dieser Ansicht. \\
\gui{Zoom selected} & Auswahl in dieser Ansicht. \\
\gui{Maximize / restore (Space)} & diese Ansicht maximieren bzw.\ zurück zum Raster. \\
\bottomrule
\end{referenz}

\begin{hinweis}
Gedreht werden kann nur in einer Ansicht vom Typ \gui{Perspective} -- auch dann, wenn sie
auf Parallelprojektion geschaltet ist. Eine Ansicht \gui{Front}, \gui{Left} oder \gui{Top}
bleibt auch mit perspektivischer Projektion in ihrer Blickrichtung. Wer zwei frei drehbare
Ansichten braucht, stellt eine zweite Ansicht auf \gui{Perspective}.
\end{hinweis}

\section{Die Konstruktionsebene}\label{sec:konstruktionsebene}
\index{Konstruktionsebene}

Jede Ansicht hat eine \emph{Konstruktionsebene}: die Ebene, in der Punkte liegen, die Sie
mit der Maus ohne Fang setzen, und auf die sich Eingaben wie \code{x,y} oder Winkel beziehen.
Sie geht immer durch den Ursprung:

\begin{referenz}{>{\raggedright\arraybackslash}p{3.2cm}>{\raggedright\arraybackslash}p{2.6cm}X}
\toprule
\textbf{Ansicht} & \textbf{Ebene} & \textbf{Achsen der Ebene (rechts / oben)} \\
\midrule
\endhead
\gui{Front} & $XZ$ & $X$ / $Z$ \\
\gui{Left} & $YZ$ & $-Y$ / $Z$ \\
\gui{Top} & $XY$ & $X$ / $Y$ \\
\gui{Perspective} & $XY$ & $X$ / $Y$ \\
\bottomrule
\end{referenz}

Die Konstruktionsebene ist als Raster dargestellt. Sein Linienabstand ist eine Zehnerpotenz
(\dots, \SI{0.1}{mm}, \SI{1}{mm}, \SI{10}{mm}, \dots), und zwar die kleinste, bei der die
Linien noch mindestens etwa 12 Pixel auseinanderliegen; jede zehnte Linie ist kräftiger.
Beim Zoomen wechselt der Abstand also in Zehnerschritten. Die beiden Achsen der Ebene sind im
Raster farbig hervorgehoben (X rot, Y grün, Z blau).

Die Koordinaten in der Statusleiste zeigen den Punkt der Konstruktionsebene unter dem
Mauszeiger -- in der Ansicht, über der sich der Zeiger gerade befindet. Rastet der Cursor auf
einen Fangpunkt ein, zeigt sie dessen Raumkoordinaten. Welche Ansicht beim Zeichnen benutzt
wird und wie sich Punkte in andere Höhen bringen lassen, beschreibt
\autoref{kap:modellieren}.

\begin{tipp}
Ein Punkt in der perspektivischen Ansicht ohne Fang liegt immer in der Ebene $Z=0$. Um auf
einer anderen Höhe zu zeichnen, benutzen Sie den Fang auf vorhandene Geometrie, tippen die
Koordinaten ein oder zeichnen in der Ansicht \gui{Front} bzw.\ \gui{Left}.
\end{tipp}

\section{Auswählen}\label{sec:auswahl}
\index{Auswahl}

Die meisten Arbeitsschritte wirken auf die \emph{Auswahl}: Löschen, Transformieren,
Gruppieren, Ausblenden, Exportieren, viele Befehle und viele Plugin-Aktionen. Ausgewählte
Objekte werden gelb dargestellt; das Dock \gui{Properties} zeigt ihre Eigenschaften, und im
Modellbaum sind ihre Zeilen markiert.

\subsection{Mit der Maus}

\begin{referenz}{>{\raggedright\arraybackslash}p{5.2cm}X}
\toprule
\textbf{Aktion} & \textbf{Wirkung} \\
\midrule
\endhead
\maus{Klick} auf ein Objekt & wählt es aus; die bisherige Auswahl wird ersetzt. \\
\maus{Klick} ins Leere & hebt die Auswahl auf. \\
\keys{Shift}+\maus{Klick} & fügt das Objekt zur Auswahl hinzu. \\
\keys{Strg}+\maus{Klick} & nimmt das Objekt aus der Auswahl heraus. \\
\keys{Strg+Shift}+\maus{Klick} & wählt eine einzelne Fläche, Kante oder einen Eckpunkt aus
  bzw.\ wieder ab (Teilobjekt, \autoref{sec:teilobjekte}). \\
\maus{links ziehen} nach rechts & Fensterauswahl: alle Objekte, die ganz im Rechteck liegen
  (\autoref{sec:rechteckauswahl}). \\
\maus{links ziehen} nach links & Kreuzen: alle Objekte, die im Rechteck liegen oder es
  berühren. \\
\bottomrule
\end{referenz}

Auf dem Mac steht \keys{Strg} für \keys{Cmd} (\autoref{sec:schreibweisen}).

\paragraph{Treffgenauigkeit.} Kurven, Kanten und Punkte werden getroffen, wenn der
Mauszeiger höchstens 6 Pixel daneben liegt. Liegen mehrere Elemente unter dem Zeiger, haben
Punkte Vorrang vor Kurven und Kanten und diese vor Flächen. In den Darstellungen
\gui{Surfaces} und \gui{Shaded} lassen sich Kurven, die hinter Flächen liegen, nicht
auswählen; schalten Sie dafür auf \gui{Wireframe}. Im Modus \gui{Wireframe} wählen Sie eine
Fläche, indem Sie eine ihrer Isolinien anklicken.

\paragraph{Klick oder Rechteck?} Bewegen Sie die Maus bei gedrückter linker Taste um mehr als
4 Pixel, wird aus dem Klick ein Auswahlrechteck. Ein leicht verwackelter Klick bleibt also
ein Klick.

\paragraph{Während eines Befehls.} Wartet ein Befehl auf einen Punkt (der Mauszeiger ist dann
ein Fadenkreuz), setzt ein Linksklick einen Punkt und wählt nichts aus. Wählen Sie deshalb
die Objekte, auf die ein Befehl wirken soll, \emph{bevor} Sie ihn starten.

\subsection{Über Menü und Tastatur}

\begin{itemize}
  \item \menu{Edit,Select All} (\keys{Strg+A}) wählt alle sichtbaren Objekte aus.
    Ausgeblendete Objekte bleiben unausgewählt.
  \item \menu{Edit,Select None / Cancel} (\keys{Esc}) hebt die Auswahl auf -- aber nur, wenn
    kein Befehl und keine Plugin-Rechnung läuft; sonst bricht \keys{Esc} zuerst diese ab.
  \item Im Modellbaum (Dock \gui{Model}) wählt ein Klick auf eine Zeile das Objekt aus, ein
    Klick auf eine Gruppe alle Objekte darin, auch in Untergruppen. \keys{Shift}- und
    \keys{Strg}-Klick erweitern die Markierung im Baum wie in einer Dateiliste
    (\autoref{sec:modellbaum}).
\end{itemize}

\subsection{Teilobjekte: Flächen, Kanten, Eckpunkte}\label{sec:teilobjekte}
\index{Teilobjekt}

Eine Hülle (\gui{Shell}) besteht aus Flächen, die an Kanten zusammenstoßen, und Eckpunkten,
an denen Kanten enden. Diese Teilobjekte lassen sich einzeln auswählen, etwa um die Fläche
einer Linse zu untersuchen, eine Kante zu vermessen oder einem Plugin eine bestimmte Fläche
anzugeben.

\begin{schritte}
  \item Halten Sie \keys{Strg+Shift} gedrückt und bewegen Sie den Mauszeiger über das Objekt.
    Statt des ganzen Objekts leuchtet jetzt die Fläche, Kante oder der Eckpunkt unter dem
    Zeiger cyan auf; die Statusleiste nennt es, z.\,B.\ \enquote{part\_1\_2 $\cdot$ Shell
    $\cdot$ edge 7}.
  \item Klicken Sie. Das Teilobjekt wird gelb und ausgewählt. Ein weiterer
    \keys{Strg+Shift}-Klick auf ein anderes Teilobjekt fügt es hinzu, ein
    \keys{Strg+Shift}-Klick auf ein ausgewähltes nimmt es wieder heraus.
  \item Das Dock \gui{Properties} zeigt bei genau einem ausgewählten Teilobjekt dessen
    Eigenschaften: bei einer Fläche Grad, Kontrollnetz, Knoten, Trimmung, Löcher und
    Flächeninhalt, bei einer Kante deren Kurvendaten, Länge und an wie vielen Flächen sie
    anliegt, bei einem Eckpunkt seine Position (\autoref{sec:eigenschaften}).
\end{schritte}

Die Nummern (\enquote{face 3}, \enquote{edge 7}, \enquote{vertex 2}) zählen die Teilobjekte
innerhalb der Hülle ab null.

\begin{hinweis}
Teilobjekte lassen sich nicht für sich löschen oder verschieben; \gui{Delete} meldet dann
\enquote{Faces, edges and vertices cannot be deleted on their own}. Befehle, die ganze Objekte
erwarten, übergehen ausgewählte Teilobjekte. \gui{Zoom selected} dagegen berücksichtigt sie.
Das Auswahlrechteck wählt immer ganze Objekte.
\end{hinweis}

\subsection{Rechteckauswahl: Fenster und Kreuzen}\label{sec:rechteckauswahl}
\index{Rechteckauswahl}\index{Fensterauswahl}\index{Kreuzen}

Ziehen Sie mit der linken Maustaste ein Rechteck auf. Die Richtung entscheidet, was ausgewählt
wird:

\begin{description}
  \item[Von links nach rechts -- Fenster.] Der Rahmen ist durchgezogen und blau. Ausgewählt
    werden nur Objekte, die \emph{vollständig} im Rechteck liegen.
  \item[Von rechts nach links -- Kreuzen.] Der Rahmen ist gestrichelt und grün. Ausgewählt
    werden alle Objekte, die im Rechteck liegen \emph{oder} es berühren.
\end{description}

Ob nach oben oder unten gezogen wird, spielt keine Rolle. Die Zusatztasten wirken wie beim
Klick: Ohne Zusatztaste ersetzt das Rechteck die Auswahl, mit \keys{Shift} fügt es hinzu, mit
\keys{Strg} nimmt es heraus.

\begin{tipp}
Um aus einer dichten Baugruppe ein inneres Teil zu greifen, ziehen Sie ein Fenster (links
nach rechts) eng um das Teil: Die umgebenden Teile ragen über das Rechteck hinaus und werden
nicht ausgewählt. Um dagegen schnell alles in einer Gegend zu erfassen, kreuzen Sie (rechts
nach links).
\end{tipp}

\subsection{Hervorhebung}\label{sec:hervorhebung}
\index{Hervorhebung}

\begin{itemize}
  \item \textbf{Cyan}: das Objekt unter dem Mauszeiger -- es würde mit einem Klick ausgewählt.
    Mit gedrückter \keys{Strg+Shift}-Taste das Teilobjekt unter dem Zeiger. Gleichzeitig nennt
    die Statusleiste Name und Art des Objekts (\enquote{Entity 12 $\cdot$ Curve}, wenn es keinen
    Namen hat).
  \item \textbf{Gelb}: ausgewählte Objekte und Teilobjekte. Flächen werden dabei gelblich
    eingefärbt, Kanten und Kurven ganz gelb gezeichnet.
\end{itemize}

Die Vorschau in Cyan zeigt also vor dem Klick, was der Klick bewirken wird -- achten Sie bei
dicht liegenden Kanten darauf, bevor Sie klicken.

\section{Ein- und Ausblenden}\label{sec:ausblenden}
\index{Ausblenden}\index{Hide}\index{Isolate}

Ausgeblendete Objekte sind unsichtbar, lassen sich nicht auswählen und nicht fangen. Sie
bleiben aber im Dokument und werden mit ihm gespeichert. Das ist praktisch, um in einer
Baugruppe an innen liegende Teile zu kommen oder um Hilfsgeometrie aus dem Weg zu räumen.

Sichtbarkeit wird an drei Stellen geschaltet; ein Objekt ist nur sichtbar, wenn alle drei es
zulassen:

\begin{enumerate}
  \item am \textbf{Objekt} selbst -- mit \gui{Hide}, \gui{Isolate} oder dem Kontrollkästchen
    in seiner Zeile im Modellbaum,
  \item an jeder \textbf{Gruppe}, in der es liegt -- Kontrollkästchen der Gruppe im
    Modellbaum (\autoref{sec:gruppen}),
  \item an seiner \textbf{Ebene} -- Kontrollkästchen im Dock \gui{Layers}
    (\autoref{sec:ebenen}).
\end{enumerate}

Die Befehle dazu stehen im Menü \menu{Edit}, in der Werkzeugleiste und im Kontextmenü des
Modellbaums:

\begin{referenz}{>{\raggedright\arraybackslash}p{3.0cm}>{\raggedright\arraybackslash}p{1.8cm}X}
\toprule
\textbf{Befehl} & \textbf{Taste} & \textbf{Wirkung} \\
\midrule
\endhead
\gui{Hide} & \keys{H} & Blendet die ausgewählten Objekte aus und hebt ihre Auswahl auf. Die
  Statusleiste meldet \enquote{\textit{n} object(s) hidden -- Unhide (Shift+H) shows them
  again}; ohne Auswahl meldet \gui{Hide} (Menü, Knopf) \enquote{Hide: select objects first}.
  Die Taste \keys{H} in einer Ansicht blendet nur mit Auswahl aus; ohne Auswahl wird sie in die
  Befehlszeile geschrieben. \\
\gui{Isolate selected} (Modellbaum: \gui{Isolate}) & -- & Blendet alle \emph{anderen} Objekte
  aus und zeigt die ausgewählten. Alle Gruppen werden dabei eingeschaltet, damit der Weg zu den
  isolierten Objekten frei ist; die Ebenen bleiben, wie sie sind. \\
\gui{Unhide} & \keys{Shift+H} & Zeigt alle Objekte wieder, die einzeln ausgeblendet waren
  (durch \gui{Hide}, \gui{Isolate} oder ihr Kontrollkästchen). Gruppen und Ebenen behalten
  ihre Schalter. Meldung \enquote{\textit{n} object(s) shown again} bzw.\ \enquote{Unhide: no
  hidden objects}. \\
\gui{Show all} & -- & Schaltet alles ein: alle Objekte, alle Gruppen und alle Ebenen. \\
\bottomrule
\end{referenz}

\paragraph{Warum hebt Hide die Auswahl auf?} Damit nichts Unsichtbares ausgewählt bleibt: Ein
anschließendes \gui{Delete} oder \gui{Move} würde sonst auf Objekte wirken, die man nicht
sieht.

\paragraph{Beispiel: an ein inneres Teil herankommen.}
\begin{schritte}
  \item Wählen Sie die Teile aus, die im Weg sind (Klick, \keys{Shift}+Klick oder
    Kreuzen-Rechteck).
  \item Drücken Sie \keys{H}. Die Teile verschwinden.
  \item Arbeiten Sie am inneren Teil.
  \item Drücken Sie \keys{Shift+H}. Die Teile sind wieder da.
\end{schritte}

\paragraph{Beispiel: nur ein Teil betrachten.}
\begin{schritte}
  \item Wählen Sie das Teil aus.
  \item \menu{Edit,Isolate selected}. Alles andere verschwindet.
  \item \keys{Z}, \keys{E} -- die Ansicht zeigt jetzt das Teil groß.
  \item \keys{Shift+H} holt die übrigen Teile zurück.
\end{schritte}

\begin{hinweis}
\gui{Unhide} und das Kontrollkästchen im Modellbaum wirken auf dieselbe Eigenschaft des
Objekts. Ist ein Objekt trotz \gui{Unhide} unsichtbar, ist seine Gruppe oder seine Ebene
ausgeschaltet -- prüfen Sie die Kontrollkästchen in \gui{Model} und \gui{Layers}, oder wählen
Sie \gui{Show all}.
\end{hinweis}

\begin{warnung}
Ein- und Ausblenden ist kein Bearbeitungsschritt: Es lässt sich nicht mit \gui{Undo}
rückgängig machen und markiert das Dokument nicht als geändert. Der Zustand wird aber beim
nächsten Speichern mitgespeichert. Ausgeblendete Objekte werden beim Export ohne Auswahl
nicht mitgeschrieben (\autoref{sec:step-export}).
\end{warnung}
